\problemname{Bygga plank}
Gustav älskar sina grannar, och har därför byggt ett plank som består av $N$ brädor. Brädorna är alla $1 \text{ dm}$ breda men kan vara olika höga. 
För att trakten ska kunna blomstra vill nu Gustav hänga upp en rektangulär anslagstavla på detta plank.
Gustav känner att anslagstavlan han hänger upp måste vara perfekt, och behöver därför din hjälp!

Mer precist vill Gustav att anslagstavlan ska vara så stor som möjligt (alltså ha så stor area som möjligt).
Anslagstavlan behöver också hänga stadigt, vilket betyder att det som mest får finnas $K$ stycken brädor bakom anslagstavlan som inte räcker hela vägen upp till dess övre kant.
För att anslagstavlan ska kunna hänga över huvud taget måste den såklart också ha minst en bräda bakom sig som räcker hela vägen upp till dess övre kant.

\section*{Indata}
Den första raden innehåller två heltal $N$ och $K$ ($0 \le K \le 20$ och $1 \le N \le 50\,000$), antalet brädor som planket består av och antalet brädor som får fattas bakom anslagstavlan.

Den andra raden innehåller $N$ heltal $h_{1}$, $h_{2}$, \dots, $h_{N}$ där $1 \le h_{i} \le 10^{9}$ är höjden på den $i$:te brädan från vänster i decimeter.

\section*{Utdata}
Skriv ut den största arean som den rektangulära anslagstavlan kan ha, i $\text{dm}^{2}$.

\section*{Poängsättning}
Din lösning kommer att testas på en mängd testfallsgrupper.
För att få poäng för en grupp så måste du klara alla testfall i gruppen.

\noindent
\begin{tabular}{| l | l | l |}
  \hline
  Grupp & Poängvärde & Gränser \\ \hline
  $1$    & $9$       &  $N \le 100, K = 1$\\ \hline 
  $2$    & $22$       &  $h_{1} < h_{2} < ... < h_{N}$ \\ \hline
  $3$    & $25$        &  $K = 0$ \\ \hline
  $4$    & $16$        &  $N \le 2000 $ \\ \hline
  $5$    & $28$       &  Inga ytterligare begränsningar. \\ \hline
\end{tabular}
